% =====================================================================
% Macros do TCC -- Cozy Pong.
% =====================================================================

% Convencao de pendencias -- ambas saem em VERMELHO e NEGRITO no PDF:
%
%   \todo{...} -> pendencia PONTUAL: uma decisao, uma referencia ou um dado
%                 que so o autor pode obter (imagem, parecer do CEP, etc.).
%
%   \ph{...}   -> PLACEHOLDER: trecho de TEXTO a redigir, tipicamente apos a
%                 coleta de dados (resultados, conclusao) ou de conteudo
%                 pessoal (agradecimentos, epigrafe).
%
% Ambas sao leves (sem tikz/todonotes); xcolor ja vem da classe JITH.
% Para localizar todas: grep -rn "\\todo{\|\\ph{" --include=*.tex .

% Robustas: aparecem dentro de \caption{}, que e argumento movel (vai para a
% Lista de Figuras/Tabelas). Sem \DeclareRobustCommand, o \color quebraria a LOF.
\DeclareRobustCommand{\todo}[1]{{\normalfont\bfseries\color{red}[TODO:\ #1]}}
\DeclareRobustCommand{\ph}[1]{{\normalfont\bfseries\color{red}[PLACEHOLDER:\ #1]}}

% Nome do jogo com estilizacao consistente (descomentar se desejar padronizar):
% \newcommand{\cozypong}{\textit{Cozy Pong}}
